\documentclass[11pt]{article}
\usepackage{amsmath,amssymb,bm}
\usepackage{geometry}
\usepackage{hyperref}
\usepackage{microtype}
\geometry{margin=1.1in}

\newcommand{\vpar}{v_{\!\parallel}}
\newcommand{\vperp}{v_{\!\perp}}
\newcommand{\vth}[1][]{v_{\mathrm{th}\ifx&#1&\else,#1\fi}}
\newcommand{\lnL}{\ln\Lambda}
\newcommand{\lambdaD}{\lambda_D}
\newcommand{\eps}{\varepsilon_0}
\newcommand{\dd}{\mathrm{d}}
\newcommand{\erf}{\operatorname{erf}}
\newcommand{\half}{\tfrac{1}{2}}

\title{Monte Carlo Coulomb Collisions for Guiding-Centre Particles\\
       in Boozer Coordinates: Implementation Notes}
\author{firm3d --- \texttt{ep\_thermal\_collisions}}
\date{\today}

\begin{document}
\maketitle
\tableofcontents
\bigskip

%-----------------------------------------------------------------------
\section{Overview}
%-----------------------------------------------------------------------

These notes describe the algorithm used in \texttt{trace\_particles\_boozer\_with\_collisions}
to include stochastic Coulomb collisions for an energetic particle (EP) tracing through
one or more Maxwellian background species.

The approach couples the deterministic guiding-centre (GC) orbit equations in Boozer
coordinates with a Milstein stochastic term for the velocity-space diffusion. The full
state vector is $(s,\theta,\zeta,v,\xi)$, where $v=|\bm{v}|$ is the total speed and
$\xi = \vpar/v$ is the pitch.

Primary references:
\begin{itemize}
  \item Hirvijoki et al.\ \cite{hirvijoki2013} --- Monte Carlo implementation of the
        guiding-centre Fokker--Planck equation in $(v,\xi)$ coordinates; Eqs.\ (32)--(33)
        of that paper define the SDE system used here.
  \item Hirvijoki et al.\ \cite{hirvijoki2014} --- ASCOT4/5 code description; source
        file \texttt{mccc\_coefs.h} defines the Coulomb logarithm convention adopted here.
  \item Boozer \& Kuo-Petravic \cite{boozer1981} --- pitch-angle representation
        in Boozer coordinates.
\end{itemize}

%-----------------------------------------------------------------------
\section{Guiding-Centre Equations in Boozer Coordinates}
%-----------------------------------------------------------------------

In Boozer coordinates $(s,\theta,\zeta)$, where $s$ is the normalised toroidal flux,
$\theta$ the poloidal angle, and $\zeta$ the toroidal angle, the vacuum GC equations are

\begin{align}
  \dot s &= -\frac{m\,v^2(1+\xi^2)}{2\,q\,\psi_0}\,\frac{\partial B}{\partial\theta}\,\frac{1}{B},
  \\[4pt]
  \dot\theta &= \frac{m\,v^2(1+\xi^2)}{2\,q\,\psi_0}\,\frac{\partial B}{\partial s}\,\frac{1}{B}
               + \frac{\iota\,\xi\,v\,B}{G},
  \\[4pt]
  \dot\zeta &= \frac{\xi\,v\,B}{G},
  \\[4pt]
  \dot\xi &= -\frac{(\iota\,\partial_\theta B + \partial_\zeta B)\,v\,(1-\xi^2)}{2G}
             \quad[\text{mirror force}],
\end{align}
where $\psi_0$ is the reference poloidal flux, $G(s)$ is the covariant toroidal
field component, and $\iota(s)$ is the rotational transform.  The magnetic moment
$\mu = v^2(1-\xi^2)/(2B)$ is reconstructed locally at each evaluation.

With collisions the equations for $v$ and $\xi$ acquire deterministic drift terms
(Section~\ref{sec:coefs}) and stochastic noise terms (Section~\ref{sec:milstein}).

%-----------------------------------------------------------------------
\section{Background Species and Profiles}
%-----------------------------------------------------------------------

Each background species $b$ is described by a \texttt{ThermalBackground} object
specifying:
\begin{itemize}
  \item $n_b(s)$ --- number density [m$^{-3}$], callable of the flux label $s\in[0,1]$;
  \item $T_b(s)$ --- temperature [J], callable (use \texttt{T\_keV $\times$ 1e3 $\times e$});
  \item $m_b$ --- species mass [kg];
  \item $q_b$ --- species charge [C] (signed).
\end{itemize}

Internally the profiles are pre-evaluated on a uniform $s$-grid of 512 points and
linearly interpolated in C++.

%-----------------------------------------------------------------------
\section{Collision Coefficients}
\label{sec:coefs}
%-----------------------------------------------------------------------

\subsection*{Notation}

The symbols used here follow the ASCOT5 source code (\texttt{mccc\_coefs.h},
\cite{ascot5code}); Table~\ref{tab:notation} shows the correspondence with
Hirvijoki et al.\ \cite{hirvijoki2013}, whose Eqs.~(32)--(33) define the
SDE structure.

\begin{table}[h]
\centering
\begin{tabular}{lll}
\hline
This paper & ASCOT5 \texttt{mccc\_coefs.h} & Hirvijoki et al.\ \cite{hirvijoki2013} \\
\hline
$D_{\parallel}$   & \texttt{Dpara}  & $D_k/2$ \\
$\nu_D$           & \texttt{nu}~$= 2D_\perp/v^2$ & $2D_\perp/v^2$ \\
$Q$               & \texttt{Q}      & (Einstein-relation drag in $K$) \\
$K$               & \texttt{K}~$= Q + D_\parallel' + 2D_\parallel/v$ & $K$ (drift in Eq.~32) \\
\hline
\end{tabular}
\caption{Notation correspondence. The noise amplitude for $v$ is
$g_v = \sqrt{2D_\parallel}$~(this paper) $= \sqrt{D_k}$~(Hirvijoki~2013).}
\label{tab:notation}
\end{table}

\subsection{Debye Length}

The total Debye screening length from all background species is \cite{spitzer1962,nrl}
\begin{equation}
  \frac{1}{\lambdaD^2}
  = \sum_b \frac{n_b\,q_b^2}{\eps\,T_b}.
  \label{eq:debye}
\end{equation}

\subsection{Coulomb Logarithm}
\label{sec:clog}

The Coulomb logarithm is defined as the ratio of maximum to minimum impact parameters
\cite{spitzer1962,nrl,ichimaru1973}:
\begin{equation}
  \lnL_{ab} = \ln\!\left(\frac{\lambdaD}{b_{\min}}\right),
  \qquad b_{\min} = \max(b_{\mathrm{cl}},\, b_{\mathrm{qm}}).
  \label{eq:lnL}
\end{equation}
The maximum impact parameter is the Debye length (Eq.~\eqref{eq:debye}).
The minimum impact parameter is the larger of the classical 90$^\circ$-deflection
distance and the quantum-mechanical de~Broglie wavelength of the reduced mass:
\begin{align}
  b_{\mathrm{cl}} &= \frac{|q_a\,q_b|}{4\pi\eps\,m_r\,\bar{v}^2},
  \label{eq:bcl}
  \\[4pt]
  b_{\mathrm{qm}} &= \frac{\hbar}{2\,m_r\,\sqrt{\bar{v}^2}},
  \label{eq:bqm}
\end{align}
where $m_r = m_a m_b/(m_a + m_b)$ is the reduced mass.

The effective velocity scale $\bar{v}^2$ is taken as
\begin{equation}
  \bar{v}^2 = v^2 + \vth[b]^2, \qquad \vth[b] = \sqrt{2T_b/m_b},
  \label{eq:vbar}
\end{equation}
following the ASCOT5 source convention \cite{ascot5code}.
This form is a numerical interpolation device, not derived from first
principles: Ichimaru \cite{ichimaru1973} derives $b_{\mathrm{cl}}$ and
$b_{\mathrm{qm}}$ only in the fast-particle limit $v\gg\vth[b]$
(Section~4.3B of that reference), where $\bar{v}\approx v$.
The $\bar{v}^2$ form extends the formula to the slow-particle regime
without a discontinuous switch:
fast EP against ions ($v\gg\vth[b]$, $\bar{v}\approx v$) and
slow EP against electrons ($v\ll\vth[e]$, $\bar{v}\approx\vth[e]$).

The quantum branch $b_{\mathrm{qm}} > b_{\mathrm{cl}}$ applies when
\begin{equation}
  \frac{2\pi\eps\,\hbar\,\sqrt{\bar{v}^2}}{|q_a\,q_b|} > 1,
\end{equation}
i.e.\ when the de~Broglie wavelength of the reduced mass particle exceeds the
classical distance of closest approach.  For 3.52\,MeV $\alpha$-particles
scattering off 10\,keV electrons one finds $b_{\mathrm{qm}} \approx 22\,b_{\mathrm{cl}}$,
so the quantum branch is active throughout alpha slowing-down against electrons.
If $\lnL \leq 0$ (Debye length below the minimum impact parameter, e.g.\
$T_b \to 0$ at finite density) the binary-collision model is undefined and
an exception is raised; ASCOT5 reaches the equivalent outcome by aborting
markers through its input-evaluation errors, and its input profiles keep
the temperature finite wherever the density is nonzero.  A warning is
printed for $0 < \lnL < 2$, where the perturbative expansion underlying
the Fokker--Planck operator begins to lose validity.

\subsection{Chandrasekhar G Function}

The Chandrasekhar function \cite{chandrasekhar1943} is
\begin{equation}
  G(x) = \frac{\erf(x) - \dfrac{2x}{\sqrt{\pi}}\,e^{-x^2}}{2x^2},
  \qquad x = v/\vth[b].
\end{equation}
Its derivative satisfies the recurrence
\begin{equation}
  G'(x) = \frac{2}{\sqrt{\pi}}\,e^{-x^2} - \frac{2\,G(x)}{x}.
\end{equation}
Key limiting behaviours:
\begin{align*}
  x\to 0:&\quad G\to 0,\quad G'\to 2/(3\sqrt{\pi})\approx 0.376; \\
  x\to\infty:&\quad G\to 1/(2x^2),\quad G'\to 0.
\end{align*}

\subsection{Collision Frequency Prefactor}

Define the Rosenbluth prefactor for species $b$ \cite{rosenbluth1957,trubnikov1965}:
\begin{equation}
  \Gamma_b = \frac{n_b\,q_a^2\,q_b^2\,\lnL_{ab}}
                  {4\pi\eps^2\,m_a^2}.
  \label{eq:Gamma}
\end{equation}

\subsection{Pitch-Angle Scattering Frequency}

\begin{equation}
  \nu_D^{(b)} = \frac{\Gamma_b\,[\erf(x) - G(x)]}{v^3}
  \;\equiv\; \frac{2D_\perp^{(b)}}{v^2}.
  \label{eq:nuD}
\end{equation}
The second form uses the perpendicular diffusion coefficient
$D_\perp^{(b)} = \Gamma_b[\erf(x)-G(x)]/(2v)$
and matches ASCOT5's \texttt{mccc\_coefs\_nu} and the $2D_\perp/v^2$
notation of Hirvijoki et al.\ \cite{hirvijoki2013,hirvijoki2014}.

\subsection{Parallel Diffusion Coefficient}

\begin{equation}
  D_{\parallel}^{(b)} = \frac{\Gamma_b\,G(x)}{v},
  \qquad
  \frac{\partial D_{\parallel}^{(b)}}{\partial v}
  = \frac{\Gamma_b}{v}\!\left(\frac{G'(x)}{\vth[b]} - \frac{G(x)}{v}\right).
  \label{eq:Dpar}
\end{equation}
In Hirvijoki et al.\ \cite{hirvijoki2013} the parallel diffusion coefficient
is denoted $D_k$, with $D_{\parallel} = D_k/2$; the noise amplitude for the
speed SDE is $\sqrt{D_k} = \sqrt{2D_{\parallel}}$ in both conventions.
ASCOT5 names this coefficient \texttt{Dpara}.
The $v$-derivative is required by the Milstein correction term
(Section~\ref{sec:milstein}).

\subsection{Deterministic Drag in $v$}

The drag coefficient (ASCOT5: \texttt{mccc\_coefs\_Q} \cite{ascot5code})
is fixed by the Einstein relation:
\begin{equation}
  Q^{(b)} = -\frac{m_a\,v}{T_b}\,D_{\parallel}^{(b)}
          = -\Gamma_b\,\frac{m_a\,G(x)}{T_b}.
  \label{eq:Q}
\end{equation}
The total deterministic drift in speed is
\begin{equation}
  K^{(b)}
  = Q^{(b)}
    + \frac{\partial D_{\parallel}^{(b)}}{\partial v}
    + \frac{2\,D_{\parallel}^{(b)}}{v}
    \label{eq:K}
\end{equation}
(ASCOT5's \texttt{mccc\_coefs\_K} \cite{ascot5code}; the drift coefficient
$K$ of Eq.~(32) in Hirvijoki et al.\ \cite{hirvijoki2013}).

\subsubsection*{It\^o transformation leading to Eq.~\eqref{eq:K}}

Equation~\eqref{eq:K} can be obtained two independent ways, which
cross-check each other.

\paragraph{(i) It\^o's formula on $v = |\bm{v}|$.}
In three-dimensional velocity space the collisional motion of the test
particle is
\begin{equation}
  \dd\bm{v} = F_{\mathrm{dyn}}\,\hat{\bm{v}}\,\dd t
            + \bm{\sigma}\cdot\dd\bm{W},
  \qquad
  \bm{\sigma}\bm{\sigma}^{\!\top}
  = 2D_{\parallel}\,\hat{\bm{v}}\hat{\bm{v}}
  + 2D_{\perp}\,(\mathsf{I}-\hat{\bm{v}}\hat{\bm{v}}),
\end{equation}
where $F_{\mathrm{dyn}}$ is the dynamical friction
\cite{chandrasekhar1943,trubnikov1965,helander2002}.  The speed
$v = |\bm{v}|$ is a nonlinear function of $\bm{v}$, so It\^o's formula
applies with $\nabla_{\!\bm v}\,v = \hat{\bm{v}}$ and Hessian
$\nabla_{\!\bm v}\nabla_{\!\bm v}\,v = (\mathsf{I}-\hat{\bm{v}}\hat{\bm{v}})/v$:
\begin{equation}
  \dd v = \Bigl(F_{\mathrm{dyn}} + \frac{2D_{\perp}}{v}\Bigr)\dd t
        + \sqrt{2D_{\parallel}}\,\dd W_v.
  \label{eq:ito_traj}
\end{equation}
The $2D_{\perp}/v$ drift is the Bessel-process effect: kicks transverse to
$\bm{v}$ do not change $v$ at first order, but they lengthen the vector at
second order, $\Delta v = |\bm{v}+\Delta\bm{v}_\perp| - v
\approx |\Delta\bm{v}_\perp|^2/2v > 0$.

\paragraph{(ii) Fokker--Planck matching (detailed balance).}
The It\^o SDE $\dd v = K\,\dd t + \sqrt{2D_{\parallel}}\,\dd W_v$ evolves
the speed density $P(v) \propto v^2 f(v)$ according to
$\partial_t P = -\partial_v(KP) + \partial_v^2(D_{\parallel}P)$.
Zero flux on the equilibrium
$P_{\mathrm{eq}} \propto v^2\,e^{-m_a v^2/2T_b}$ requires
$K P_{\mathrm{eq}} = \partial_v(D_{\parallel} P_{\mathrm{eq}})$, i.e.
\begin{equation}
  K = -\frac{m_a v}{T_b}\,D_{\parallel}
    + \frac{\partial D_{\parallel}}{\partial v}
    + \frac{2\,D_{\parallel}}{v},
\end{equation}
which is Eq.~\eqref{eq:K} with $Q$ as in Eq.~\eqref{eq:Q}.  Here
$\partial_v D_{\parallel}$ is the multiplicative-noise (It\^o) drift that
appears when the diffusion term is commuted into advective form, and
$2D_{\parallel}/v$ comes from the Jacobian $v^2$ of spherical
velocity-space coordinates.

\paragraph{Consistency of (i) and (ii).}
For Maxwellian field particles the dynamical friction obeys the identity
\begin{equation}
  F_{\mathrm{dyn}}^{(b)}
  = -\Bigl(1+\frac{m_a}{m_b}\Bigr)\frac{2x^2\,G(x)\,\Gamma_b}{v^2}
  = Q^{(b)} + \frac{\partial D_{\parallel}^{(b)}}{\partial v}
    + \frac{2\,\bigl(D_{\parallel}^{(b)}-D_{\perp}^{(b)}\bigr)}{v},
  \label{eq:fdyn}
\end{equation}
so that $F_{\mathrm{dyn}} + 2D_{\perp}/v
= Q + \partial_v D_{\parallel} + 2D_{\parallel}/v = K$ exactly.
($F_{\mathrm{dyn}}$ is ASCOT5's \texttt{mccc\_coefs\_F}, which is not used
in the guiding-centre scheme.)  In the fast-particle limit $x \gg 1$ this
identity encodes the classic cancellation
\begin{equation*}
  F_{\mathrm{dyn}} + \frac{2D_{\perp}}{v}
  \;\longrightarrow\;
  -\Bigl(1+\frac{m_a}{m_b}\Bigr)\frac{\Gamma_b}{v^2}
  + \frac{\Gamma_b}{v^2}
  = -\frac{m_a}{m_b}\,\frac{\Gamma_b}{v^2}:
\end{equation*}
most of the dynamical friction is compensated by transverse diffusion, and
the net slowing-down rate is set by the mass ratio.  Note in particular
that $Q$ is \emph{not} the dynamical friction --- it is the
Einstein-relation drag $-(m_a v/T_b)D_{\parallel}$.

\subsection{Summation over Species}

All coefficients are additive over background species:
\begin{equation}
  \nu_D = \sum_b \nu_D^{(b)}, \quad
  D_{\parallel} = \sum_b D_{\parallel}^{(b)}, \quad
  K = \sum_b K^{(b)}.
\end{equation}
The Debye length used in Eq.~\eqref{eq:lnL} is computed from all species
simultaneously (Eq.~\eqref{eq:debye}), prior to the per-species loop.

%-----------------------------------------------------------------------
\section{Stochastic Differential Equations}
%-----------------------------------------------------------------------

\subsection{Full SDE System}

The complete SDE for the five-dimensional state $(s,\theta,\zeta,v,\xi)$ is
\cite{hirvijoki2013}

\begin{align}
  \dd s      &= \dot s_{\mathrm{orb}}\,\dd t, \\
  \dd\theta  &= \dot\theta_{\mathrm{orb}}\,\dd t, \\
  \dd\zeta   &= \dot\zeta_{\mathrm{orb}}\,\dd t, \\
  \dd v      &= K(v,s)\,\dd t
               + \sqrt{2D_{\parallel}}\,\dd W_v, \label{eq:sde_v}\\
  \dd\xi     &= \dot\xi_{\mathrm{mirror}}\,\dd t
               - \xi\,\nu_D\,\dd t
               + \sqrt{\nu_D(1-\xi^2)}\,\dd W_\xi, \label{eq:sde_xi}
\end{align}
where $\dd W_v$ and $\dd W_\xi$ are independent Wiener increments with
$\mathbb{E}[\dd W^2]=\dd t$.
Equations~\eqref{eq:sde_v}--\eqref{eq:sde_xi} correspond to
Eqs.\ (32)--(33) of Hirvijoki et al.\ \cite{hirvijoki2013}, with the
identification $D_\parallel \leftrightarrow D_k/2$ and $\nu_D \leftrightarrow 2D_\perp/v^2$
in their notation.  The noise amplitude for $\xi$ ensures that the marginal
distribution of $\xi$ converges to a uniform distribution on $[-1,1]$ in the
isotropisation limit $\nu_D t\gg 1$.

\subsection{Milstein Scheme}
\label{sec:milstein}

After each accepted Dormand--Prince (DP) step of physical size $h$, the stochastic
increments $\Delta W_v\sim\mathcal{N}(0,h)$ and $\Delta W_\xi\sim\mathcal{N}(0,h)$
are generated and applied via the Milstein method \cite{kloeden1992}:

\begin{align}
  v  &\leftarrow v
     + g_v\,\Delta W_v
     + \half\frac{\partial g_v}{\partial v}\bigl(\Delta W_v^2 - h\bigr),
  \\[4pt]
  \xi &\leftarrow \xi
     + g_\xi\,\Delta W_\xi
     + \half\frac{\partial g_\xi}{\partial \xi}\bigl(\Delta W_\xi^2 - h\bigr),
\end{align}
where the noise amplitudes and their derivatives are
\begin{align}
  g_v  &= \sqrt{2D_{\parallel}},
  &\frac{\partial g_v}{\partial v} &= \frac{\partial D_{\parallel}/\partial v}{\sqrt{2D_{\parallel}}},
  \\[4pt]
  g_\xi &= \sqrt{\nu_D(1-\xi^2)},
  &\frac{\partial g_\xi}{\partial \xi} &= \frac{-\xi\,\nu_D}{\sqrt{\nu_D(1-\xi^2)}}.
\end{align}
In terms of the coefficient structs, the Milstein correction for $\xi$ simplifies to
\begin{equation}
  \delta\xi_{\mathrm{Milstein}}
  = -\half\,\xi\,\nu_D\,\bigl(\Delta W_\xi^2 - h\bigr).
\end{equation}
Boundary conditions follow ASCOT5 (\texttt{mccc\_gc\_milstein.c},
\cite{ascot5code}).  The speed reflects off a thermal cutoff,
\begin{equation}
  v < v_{\mathrm{cut}} \;\Rightarrow\; v \leftarrow 2\,v_{\mathrm{cut}} - v,
  \qquad
  v_{\mathrm{cut}} = 0.1\,\sqrt{T_{\min}/m_a},
\end{equation}
where $T_{\min}$ is the coldest active background species (ASCOT5's
\texttt{MCCC\_CUTOFF} uses its first background species), so that the
collision coefficients, which diverge as $v\to0$, are never evaluated deep
below the thermal speed.  The pitch mirrors at the boundary:
$|\xi| > 1 \Rightarrow \xi \leftarrow \mathrm{sign}(\xi)\,(2 - |\xi|)$;
a hard clamp would pile up probability at exactly $\xi = \pm 1$.

%-----------------------------------------------------------------------
\section{Adaptive Integration Strategy}
%-----------------------------------------------------------------------

\subsection{Dormand--Prince Solver}

The deterministic ODE part (all five components of the drift) is integrated
with the classic 4(5)-order embedded Dormand--Prince (FSAL) method \cite{dormand1980}.
Step-size control uses mixed absolute/relative tolerances:
\begin{equation}
  \mathrm{err}_i = \bigl|y_i^{(5)} - y_i^{(4)}\bigr|
  \leq \delta_{\mathrm{abs}} + \delta_{\mathrm{rel}}\,|y_i|.
\end{equation}
Default tolerances: $\delta_{\mathrm{abs}}=\delta_{\mathrm{rel}}=10^{-9}$.

\subsection{Minimum Step Size}

The parameter \texttt{DP\_hmin} (default \texttt{0.0}) sets a lower bound on the
physical step size $h$ in seconds.  A value of $10^{-10}$--$10^{-9}$\,s prevents the
solver from stalling at bounce points while keeping steps small enough that the orbital
dynamics remain accurate.

\subsection{Operator Splitting}

The SDE is solved via Lie--Trotter operator splitting:
\begin{enumerate}
  \item \textbf{Deterministic half:} advance $(s,\theta,\zeta,v,\xi)$ by one DP step
        of physical time $h$ using the full drift
        $(\dot s,\dot\theta,\dot\zeta,K,-\xi\nu_D+\dot\xi_{\mathrm{mirror}})$.
  \item \textbf{Stochastic half:} apply Milstein noise to $(v,\xi)$ at the accepted
        state, using $h$ as the diffusion window.
\end{enumerate}
The trajectory is saved \emph{before} the Milstein step at each
$t_{\mathrm{save}}$ checkpoint, so saved states correspond to the deterministic
prediction (consistent with the standard output convention).

%-----------------------------------------------------------------------
\section{Output Format}
%-----------------------------------------------------------------------

Each particle trace is a NumPy array of shape $(\text{ntimesteps},\,6)$ with columns
\[
  [t,\; s,\; \theta,\; \zeta,\; \vpar,\; v].
\]
From these the kinetic energy $E = \half m v^2$ and the magnetic moment
$\mu = (v^2 - \vpar^2)/(2B)$ can be reconstructed at each saved point.
When \texttt{forget\_exact\_path=True} only the first and last rows are retained.

Each particle uses RNG seed $\texttt{rng\_seed} + i$, making multi-particle
MPI runs fully reproducible.

%-----------------------------------------------------------------------
\section{Python Interface}
%-----------------------------------------------------------------------

\begin{verbatim}
from firm3d.field.collisions import ThermalBackground, \
    trace_particles_boozer_with_collisions
from firm3d.util.constants import (
    PROTON_MASS, ELECTRON_MASS, ELEMENTARY_CHARGE, ONE_EV,
    ALPHA_PARTICLE_MASS, ALPHA_PARTICLE_CHARGE,
    FUSION_ALPHA_PARTICLE_ENERGY,
)

# Typical DT-fusion plasma backgrounds
n0 = 1e20          # m^-3
deuterium = ThermalBackground(
    n_profile=lambda s: n0 * (1 - s),
    T_profile=lambda s: 10e3 * ONE_EV * (1 - s),
    mass=2 * PROTON_MASS,
    charge=ELEMENTARY_CHARGE,
)
tritium = ThermalBackground(
    n_profile=lambda s: n0 * (1 - s),
    T_profile=lambda s: 10e3 * ONE_EV * (1 - s),
    mass=3 * PROTON_MASS,
    charge=ELEMENTARY_CHARGE,
)
electrons = ThermalBackground(
    n_profile=lambda s: 2 * n0 * (1 - s),   # quasi-neutral
    T_profile=lambda s: 10e3 * ONE_EV * (1 - s),
    mass=ELECTRON_MASS,
    charge=-ELEMENTARY_CHARGE,
)

res_tys, res_hits = trace_particles_boozer_with_collisions(
    field,
    stz_inits,                              # (N, 3) array
    parallel_speeds,                        # (N,) array, m/s
    backgrounds=[deuterium, tritium, electrons],
    tmax=1e-1,                              # 100 ms
    mass=ALPHA_PARTICLE_MASS,
    charge=ALPHA_PARTICLE_CHARGE,
    Ekin=FUSION_ALPHA_PARTICLE_ENERGY,      # 3.52 MeV
    tol=1e-9,
    dt_save=1e-4,                           # save every 0.1 ms
    DP_hmin=1e-10,                          # s
    rng_seed=42,
)
\end{verbatim}

%-----------------------------------------------------------------------
\section{Validity and Limitations}
%-----------------------------------------------------------------------

\begin{itemize}
  \item \textbf{Vacuum field:} electric fields, $E\times B$ drifts, and
        electromagnetic perturbations are not included in the collision model
        (though perturbations can be combined via the existing SAW tracing infrastructure).

  \item \textbf{Maxwellian backgrounds:} the collision operator assumes each
        background species has a local Maxwellian distribution function.  Distorted
        distributions (e.g.\ beam-driven) are not supported.

  \item \textbf{Test-particle limit:} the EP density is assumed negligibly small
        compared to the background, so no back-reaction on the background is included.

  \item \textbf{Guiding-centre approximation:} valid when the Larmor radius
        $r_L = mv_\perp/(|q|B)$ is small compared to the scale lengths of $B$ and the
        density/temperature profiles.  For 3.52\,MeV $\alpha$-particles in a
        $5$\,T stellarator $r_L\sim 5$\,cm, comparable to device scale, so
        finite-Larmor-radius corrections may be non-negligible.

  \item \textbf{Zeroth-order GC collision operator:} the implementation follows
        Hirvijoki et al.\ \cite{hirvijoki2013} at zeroth order in $\rho_*/L$
        (Larmor radius over gradient length).  The spatial diffusion term
        $D_X \propto D_\perp/\Omega_g^2$ present in the full operator
        \cite{hirvijoki2014} is neglected.

  \item \textbf{Slow tests: \texttt{@pytest.mark.slow}} --- five physics-validation
        tests require $N\sim 20$ particles integrated for $5$--$100\,\mu$s and are
        intended for HPC resources (Perlmutter or similar).
\end{itemize}

%-----------------------------------------------------------------------
\begin{thebibliography}{99}
%-----------------------------------------------------------------------

\bibitem{hirvijoki2013}
E.~Hirvijoki, A.~Brizard, A.~Snicker, and T.~Kurki-Suonio,
``Monte Carlo implementation of a guiding-center Fokker-Planck kinetic
equation,''
\textit{Phys.\ Plasmas} \textbf{20}, 092505 (2013).
\texttt{doi:10.1063/1.4820951}

\bibitem{hirvijoki2014}
E.~Hirvijoki et al.,
``ASCOT: Solving the kinetic equation of minority particle species in
tokamak plasmas,''
\textit{Comput.\ Phys.\ Commun.}\ \textbf{185}, 1310 (2014).
\texttt{doi:10.1016/j.cpc.2013.10.019}

\bibitem{ascot5code}
ASCOT5 source code, \texttt{src/simulate/mccc/mccc\_coefs.h},
function \texttt{mccc\_coefs\_clog},
\url{https://github.com/ascot4fusion/ascot5/blob/main/src/simulate/mccc/mccc_coefs.h}.

\bibitem{boozer1981}
A.~H.~Boozer and G.~Kuo-Petravic,
``Monte Carlo evaluation of transport coefficients,''
\textit{Phys.\ Fluids} \textbf{24}, 851 (1981).
\texttt{doi:10.1063/1.863445}

\bibitem{spitzer1962}
L.~Spitzer,
\textit{Physics of Fully Ionized Gases}, 2nd ed.\
(Interscience, New York, 1962), Chapter~5.

\bibitem{nrl}
J.~D.~Huba,
\textit{NRL Plasma Formulary}
(Naval Research Laboratory, Washington DC, 2019).
Available at \url{https://www.nrl.navy.mil/ppd/content/nrl-plasma-formulary}

\bibitem{ichimaru1973}
S.~Ichimaru,
\textit{Basic Principles of Plasma Physics: A Statistical Approach}
(Benjamin, Reading MA, 1973).

\bibitem{chandrasekhar1943}
S.~Chandrasekhar,
``Dynamical friction I: General considerations,''
\textit{Astrophys.\ J.} \textbf{97}, 255 (1943).
\texttt{doi:10.1086/144517}

\bibitem{rosenbluth1957}
M.~N.~Rosenbluth, W.~M.~MacDonald, and D.~L.~Judd,
``Fokker-Planck equation for an inverse-square force,''
\textit{Phys.\ Rev.} \textbf{107}, 1 (1957).
\texttt{doi:10.1103/PhysRev.107.1}

\bibitem{trubnikov1965}
B.~A.~Trubnikov,
``Particle interactions in a fully ionized plasma,''
in \textit{Reviews of Plasma Physics}, Vol.~1,
M.~A.~Leontovich, ed.\ (Consultants Bureau, New York, 1965), p.~105.

\bibitem{helander2002}
P.~Helander and D.~J.~Sigmar,
\textit{Collisional Transport in Magnetized Plasmas}
(Cambridge University Press, Cambridge, 2002), Chapter~3.

\bibitem{kloeden1992}
P.~E.~Kloeden and E.~Platen,
\textit{Numerical Solution of Stochastic Differential Equations}
(Springer, Berlin, 1992).

\bibitem{dormand1980}
J.~R.~Dormand and P.~J.~Prince,
``A family of embedded Runge-Kutta formulae,''
\textit{J.\ Comput.\ Appl.\ Math.} \textbf{6}, 19 (1980).
\texttt{doi:10.1016/0771-050X(80)90013-3}

\end{thebibliography}

\end{document}
